\section{Interoperability}\label{sec:interop}

A GPU-native logic engine is only as useful as its ability to exchange data with the frameworks researchers already use. xlog exposes GPU-resident relations, probabilities, and gradients through four standard interfaces---DLPack, Arrow, Python bindings, and PyTorch autograd---so that results flow into downstream tools without device-to-host copies or format conversion.

\subsection{DLPack: Zero-Copy GPU Tensor Sharing}

xlog implements the DLPack~\cite{dlpack} protocol for bidirectional zero-copy exchange of GPU buffers with any framework that speaks it---PyTorch, JAX, CuPy, cuDF. On export, each relation column becomes a managed DLPack tensor whose underlying GPU memory is reference-counted and freed only once every consumer releases its handle. On import, xlog validates device affinity, dtype, contiguity, and alignment---optionally against an expected schema, rejecting mismatches before evaluation---and assembles the tensors into a device buffer with no host copy. All eight scalar types have well-defined DLPack dtype mappings. In Python the protocol surfaces as a PyCapsule, so a consumer calls \texttt{torch.from\_dlpack(capsule)} to obtain a live CUDA tensor backed by the very memory xlog wrote.

\subsection{Arrow IPC: Columnar Export with Dictionary Encoding}

For columnar analytics tools---Pandas, Polars, DuckDB---xlog exports relations through the Arrow~\cite{arrow} C Device Data Interface, pairing Arrow array and schema pointers with a CUDA device type and ordinal for zero-copy round-tripping. Symbol-typed columns get special handling: interned symbol IDs export as an Arrow dictionary array over a string dictionary, so a Polars or DuckDB consumer reads symbolic columns as native string dictionaries with no separate lookup, while xlog keeps the compact integer representation internally.

\subsection{Python Bindings}

The \texttt{pyxlog} crate provides Python bindings built with PyO3~\cite{pyo3}, exposing compilation entry points for deterministic, probabilistic, and ILP programs, each returning a compiled object with evaluation, training, and result-extraction methods; type stubs give full IDE auto-completion and static checking. Tensor-valued results are returned as DLPack capsules, keeping data on the GPU and letting the caller choose how to materialize it. Long-running GPU operations release the Python GIL before kernel dispatch---circuit evaluation, ILP fixpoint iteration, and the neural backward pass all run without holding it, so Python data-loader threads and asyncio tasks proceed concurrently with GPU computation.

\subsection{PyTorch Autograd Integration}

Neural-symbolic training requires gradients to flow from logical inference back into neural network parameters. xlog achieves this by constructing loss tensors that participate directly in PyTorch's autograd graph: it evaluates the compiled circuit on the GPU---with the GIL released---then exports the scalar NLL loss as a DLPack capsule and imports it via \texttt{torch.from\_dlpack}. Because the network's forward pass produced grad-enabled tensors that fed into the circuit, the loss tensor stays connected to the autograd graph. A batched path stacks per-network inputs into one forward pass, evaluates multiple queries against the shared circuit, accumulates per-query losses, and runs a single \texttt{torch.autograd.backward}. The circuit thus behaves as a differentiable function inside PyTorch---gradients flow from the logic layer through DLPack tensors back to \texttt{nn.Module} parameters, with no custom autograd \texttt{Function} subclass required.
